% Supersonic technical drawing
% Author: Michele Muccioli
% Modified Version: ADA-DRAW, Fabian Schöfer

\documentclass{standalone}
\usepackage{tikz}
\usetikzlibrary{calc,patterns,arrows,shapes.arrows,intersections}
\usetikzlibrary{decorations}
\usepackage{wasysym}
\usepackage{siunitx}
\usepackage{xfp} % in the preamble
%%%%%%%%%%%%%%%%%%%%%%%%%%%%%%
% SECTION PATTERN DEFINITION %
%%%%%%%%%%%%%%%%%%%%%%%%%%%%%%
\newlength\thicknes
\pgfdeclarepatternformonly[\thickness]{section}
{\pgfpointorigin}
{\pgfpoint{1cm}{1cm}}
{\pgfpoint{1cm}{1cm}}
{
\pgfsetlinewidth{\thickness}
\pgfpathmoveto{\pgfpoint{0cm}{0cm}}
\pgfpathlineto{\pgfpoint{1cm}{1cm}}
\pgfpathclose
\pgfsetlinewidth{\thickness}
\pgfpathmoveto{\pgfpoint{0cm}{.5cm}}
\pgfpathlineto{\pgfpoint{.5cm}{1cm}}
\pgfpathclose
\pgfsetlinewidth{\thickness}
\pgfpathmoveto{\pgfpoint{.5cm}{0cm}}
\pgfpathlineto{\pgfpoint{1cm}{.5cm}}
\pgfusepath{stroke}
}

\tikzset{
thickness/.store in = \thickness,
thickness           = 0.5pt
}
%%%%%%%%%%%%%%%%%%%%%%%%%
% DIMENSION DECORATIONS %
%%%%%%%%%%%%%%%%%%%%%%%%%
\makeatletter
%-> New if
\newif\if@dim@connection

%-> New TeX dimensions
\newdimen\dim@x
\newdimen\dim@y
\newdimen\dim@sep
\newdimen\dim@overline
\newdimen\dim@overlineI
\newdimen\dim@overlineII
\newdimen\dim@text@translation


%-> Horizontal dimension
\pgfdeclaredecoration{Hdim}{final}{%
\state{final}{%
% Setting needed dimensions from pgfkeys values.
\tikz@checkunit{\pgfkeysvalueof{/pgf/decoration/distance}}
\iftikz@isdimension
	\pgfmathsetlength{\dim@sep}{\pgfkeysvalueof{/pgf/decoration/distance}}
\else
	\pgfmathsetlength{\dim@sep}{\pgfkeysvalueof{/pgf/decoration/distance}*\pgf@yy}
\fi
\pgfmathsetlength{\dim@x}{\pgfdecoratedpathlength*cos(\pgfdecoratedangle)}
\pgfmathsetlength{\dim@y}{\pgfdecoratedpathlength*sin(\pgfdecoratedangle)}
\pgfmathsetlength{\dim@overline}{\pgfkeysvalueof{/pgf/decoration/overline}}
\pgfmathsetlength{\dim@text@translation}{\pgfkeysvalueof{/pgf/decoration/text translation}}
% Setting text to write from pgfkeys value.
\def\dim@text{\pgfkeysvalueof{/pgf/decoration/text}}
% The overline verse (where it points. Positive or negative) depends
% on dimension sep and angle of the segment to dimension.
\ifdim0pt<\dim@sep
	% Case with positive dimension sep.
	% Normal behaviour.
	\pgfmathsetlength{\dim@overlineI}{\dim@overline}
	\pgfmathsetlength{\dim@overlineII}{\dim@overline}
\else
	% Case with negative dimension sep.
	% The value of dimension overline I is different from
	% dimension overline II and it depends on the
	% quadrant where the the segment is.
	\pgfmathparse{abs(\dim@sep)<abs(\dim@y)}
	\if\pgfmathresult1
		% true
		\ifnum180<\pgfdecoratedangle
			\pgfmathsetlength{\dim@overlineI}{-\dim@overline}
			\pgfmathsetlength{\dim@overlineII}{\dim@overline}
		\else
			\pgfmathsetlength{\dim@overlineI}{\dim@overline}
			\pgfmathsetlength{\dim@overlineII}{-\dim@overline}
		\fi
	\else
		% false
		\pgfmathsetlength{\dim@overlineI}{-\dim@overline}
		\pgfmathsetlength{\dim@overlineII}{-\dim@overline}
	\fi
\fi
% Set to rotation transformation because we want dimension
% only the horizontal distance between two points.
\pgftransformrotate{-\pgfdecoratedangle}
% The drawing depends on segment to dimension orientation
\ifnum180<\pgfdecoratedangle
	% FIRST POINT HIGHER
	%
	% Node drawing. It can be translate along the path with
	% text translation option.
	{
 	\pgftransformshift{\pgfpoint{\dim@x/2+\dim@text@translation}{\dim@sep}}
	% Capturing the TikZ picture font size
	\dim@text@font
	\pgfnode{rectangle}{south}{\dim@text}{\dim@text@nodename}{\pgfusepath{discard}}
	}
	% Drawing the connection segments.
	\pgfpathmoveto{\pgfpoint{0pt}{\dim@sep+\dim@overlineI}}
	% Check if the first connection line must be drawn.
	\if@dim@connection
		\pgfpathlineto{\pgfpoint{0pt}{0pt}}
	\else
		\pgfpathlineto{\pgfpoint{0pt}{\dim@sep-\dim@overlineI}}
	\fi 
	\pgfpathmoveto{\pgfpoint{\dim@x}{\dim@sep+\dim@overlineII}}
	% Check if the second connection line must be drawn.
	\if@dim@connection
		\pgfpathlineto{\pgfpoint{\dim@x}{\dim@y}}
	\else
		\pgfpathlineto{\pgfpoint{\dim@x}{\dim@sep-\dim@overlineII}}
	\fi
	% Draw an extra line if node text lies outside of the
	% dimension.
	\pgfmathparse{abs(.5\dim@x)<abs(\dim@text@translation)}
	\if\pgfmathresult1
		% Dimension label outside.
		\ifdim0pt<\dim@text@translation
			% Text lies on right.
			\pgfpathmoveto{\pgfpointanchor{\dim@text@nodename}{south east}}
			\pgfpathlineto{\pgfpoint{\dim@x}{\dim@sep}}
		\else
			% Text lies on left.
			\pgfpathmoveto{\pgfpointanchor{\dim@text@nodename}{south west}}
			\pgfpathlineto{\pgfpoint{0pt}{\dim@sep}}
		\fi
	\fi
	% Draw the dimension segments with arrows.
	\pgfsetarrowsstart{\dim@arrow@type}
	\pgfsetarrowsend{\dim@arrow@type}
	\pgfpathmoveto{\pgfpoint{0pt}{\dim@sep}}
	\pgfpathlineto{\pgfpoint{\dim@x}{\dim@sep}}
	% Finally put all the stuff on paper.
	\pgfusepath{stroke}
\else
	% FIRST POINT LOWER
	%
	% Node drawing. It can be translate along the path with
	% text translation option.
	{
 	\pgftransformshift{\pgfpoint{\dim@x/2+\dim@text@translation}{\dim@y+\dim@sep}}
	% Capturing the TikZ picture font size
	\dim@text@font
	\pgfnode{rectangle}{south}{\dim@text}{\dim@text@nodename}{\pgfusepath{discard}}
	}
	% Drawing the connection segments.
	\pgfpathmoveto{\pgfpoint{0pt}{\dim@y+\dim@sep+\dim@overlineI}}
	% Check if the first connection line must be drawn.
	\if@dim@connection
		\pgfpathlineto{\pgfpoint{0pt}{0pt}}
	\else
		\pgfpathlineto{\pgfpoint{0pt}{\dim@y+\dim@sep-\dim@overlineI}}		
	\fi
	\pgfpathmoveto{\pgfpoint{\dim@x}{\dim@y+\dim@sep+\dim@overlineII}}
	% Check if the second connection line must be drawn.
	\if@dim@connection
		\pgfpathlineto{\pgfpoint{\dim@x}{\dim@y}}
	\else
		\pgflineto{\pgfpoint{\dim@x}{\dim@y+\dim@sep-\dim@overlineII}}
	\fi
	% Draw an extra line if node text lies outside of the
	% dimension.
	\pgfmathparse{abs(.5\dim@x)<abs(\dim@text@translation)}
	\if\pgfmathresult1
		% Dimension label outside.
		\ifdim0pt<\dim@text@translation
			% Text lies on right.
			\pgfpathmoveto{\pgfpointanchor{\dim@text@nodename}{south east}}
			\pgfpathlineto{\pgfpoint{\dim@x}{\dim@y+\dim@sep}}
		\else
			% Text lies on left.
			\pgfpathmoveto{\pgfpointanchor{\dim@text@nodename}{south west}}
			\pgfpathlineto{\pgfpoint{0pt}{\dim@y+\dim@sep}}
		\fi
	\fi
	% Draw the dimension segments with arrows.
	\pgfsetarrowsstart{\dim@arrow@type}
	\pgfsetarrowsend{\dim@arrow@type}
	\pgfpathmoveto{\pgfpoint{0pt}{\dim@y+\dim@sep}}
	\pgfpathlineto{\pgfpoint{\dim@x}{\dim@y+\dim@sep}}
	% Finally put all the stuff on paper.
	\pgfusepath{stroke}
\fi
}}


%-> Vertical dimension (NOT UTILISED IN THE FOLLOWING DRAWING)
\pgfdeclaredecoration{Vdim}{final}{%
\state{final}{%
% Setting needed dimensions from pgfkeys values.
\tikz@checkunit{\pgfkeysvalueof{/pgf/decoration/distance}}
\iftikz@isdimension
	\pgfmathsetlength{\dim@sep}{\pgfkeysvalueof{/pgf/decoration/distance}}
\else
	\pgfmathsetlength{\dim@sep}{\pgfkeysvalueof{/pgf/decoration/distance}*\pgf@yy}
\fi
\pgfmathsetlength{\dim@x}{\pgfdecoratedpathlength*cos(\pgfdecoratedangle)}
\pgfmathsetlength{\dim@y}{\pgfdecoratedpathlength*sin(\pgfdecoratedangle)}
\pgfmathsetlength{\dim@overline}{\pgfkeysvalueof{/pgf/decoration/overline}}
\pgfmathsetlength{\dim@text@translation}{\pgfkeysvalueof{/pgf/decoration/text translation}}
% Setting text to write from pgfkeys value.
\def\dim@text{\pgfkeysvalueof{/pgf/decoration/text}}
% The overline verse (where it points. Positive or negative) depends
% on dimension sep and angle of the segment to dimension.
\ifdim0pt<\dim@sep
	% Case with positive dimension sep.
	% Normal behaviour.
	\pgfmathsetlength{\dim@overlineI}{\dim@overline}
	\pgfmathsetlength{\dim@overlineII}{\dim@overline}
\else
	% Case with negative dimension sep.
	% The value of dimension overline I is different from
	% dimension overline II and it depends on the
	% quadrant where the the segment is.
	\pgfmathparse{abs(\dim@sep)<abs(\dim@x)}
	\if\pgfmathresult1
		% true
		\ifnum90<\pgfdecoratedangle
			\ifnum270<\pgfdecoratedangle
				\pgfmathsetlength{\dim@overlineI}{\dim@overline}
				\pgfmathsetlength{\dim@overlineII}{-\dim@overline}
			\else
				\pgfmathsetlength{\dim@overlineI}{-\dim@overline}
				\pgfmathsetlength{\dim@overlineII}{\dim@overline}
			\fi
		\else
			\pgfmathsetlength{\dim@overlineI}{\dim@overline}
			\pgfmathsetlength{\dim@overlineII}{-\dim@overline}
		\fi
	\else
		% false
		\pgfmathsetlength{\dim@overlineI}{-\dim@overline}
		\pgfmathsetlength{\dim@overlineII}{-\dim@overline}
	\fi
\fi
% Set to rotation transformation because we want dimension
% only the horizontal distance between two points.
\pgftransformrotate{-\pgfdecoratedangle}
\ifnum180<\pgfdecoratedangle
	%
	% FIRST POINT HIGHER
	%
	\ifnum270>\pgfdecoratedangle
		% Third quadrant
		% Node drawing. It can be translate along the path with
		% text translation option.
		{
 		\pgftransformshift{\pgfpoint{\dim@sep}{\dim@y/2-\dim@text@translation}}
		% In this case, the node must e rotated.
		\pgftransformrotate{-90}
		% Capturing the TikZ picture font size
		\dim@text@font
		\pgfnode{rectangle}{south}{\dim@text}{\dim@text@nodename}{\pgfusepath{discard}}
		}
		% Drawing the connection segments.
		\pgfpathmoveto{\pgfpoint{\dim@sep+\dim@overlineI}{0pt}}
		% Check if the first connection line must be drawn.
		\if@dim@connection
			\pgfpathlineto{\pgfpoint{0pt}{0pt}}
		\else
			\pgfpathlineto{\pgfpoint{\dim@sep-\dim@overlineI}{0pt}}
		\fi
		\pgfpathmoveto{\pgfpoint{\dim@sep+\dim@overlineII}{\dim@y}}
		% Check if the second connection line must be drawn.
		\if@dim@connection
			\pgfpathlineto{\pgfpoint{\dim@x}{\dim@y}}
		\else
			\pgfpathlineto{\pgfpoint{\dim@sep-\dim@overlineII}{\dim@y}}
		\fi
		% Draw an extra line if node text lies outside of the
		% dimension.
		\pgfmathparse{abs(.5\dim@y)<abs(\dim@text@translation)}
		\if\pgfmathresult1
			% Dimension label outside.
			\ifdim0pt<\dim@text@translation
				% Text lies on right.
				\pgfpathmoveto{\pgfpointanchor{\dim@text@nodename}{south east}}
				\pgfpathlineto{\pgfpoint{\dim@sep}{\dim@y}}
			\else
				% Text lies on left.
				\pgfpathmoveto{\pgfpointanchor{\dim@text@nodename}{south west}}
				\pgfpathlineto{\pgfpoint{\dim@sep}{0pt}}
			\fi
		\fi
		% Draw the dimension segments with arrows. 
		\pgfsetarrowsstart{\dim@arrow@type}
		\pgfsetarrowsend{\dim@arrow@type}
		\pgfpathmoveto{\pgfpoint{\dim@sep}{0pt}}
		\pgfpathlineto{\pgfpoint{\dim@sep}{\dim@y}}
		% Finally put all the stuff on paper.
		\pgfusepath{stroke}
	\else
		% Fourth quadrant
		% Node drawing. It can be translate along the path with
		% text translation option.
		{
 		\pgftransformshift{\pgfpoint{\dim@x+\dim@sep}{\dim@y/2-\dim@text@translation}}
		% In this case, the node must e rotated.
		\pgftransformrotate{-90}
		% Capturing the TikZ picture font size
		\dim@text@font
		\pgfnode{rectangle}{south}{\dim@text}{\dim@text@nodename}{\pgfusepath{discard}}
		}
		% Drawing the connection segments.
		\pgfpathmoveto{\pgfpoint{\dim@x+\dim@sep+\dim@overlineI}{0pt}}
		% Check if the first dimension connection must be drawn.
		\if@dim@connection
			\pgfpathlineto{\pgfpoint{0pt}{0pt}}
		\else
			\pgfpathlineto{\pgfpoint{\dim@x+\dim@sep-\dim@overlineI}{0pt}}
		\fi
		\pgfpathmoveto{\pgfpoint{\dim@x+\dim@sep+\dim@overlineII}{\dim@y}}
		% Check if the second dimension connection must be drawn.
		\if@dim@connection
			\pgfpathlineto{\pgfpoint{\dim@x}{\dim@y}}
		\else
			\pgfpathlineto{\pgfpoint{\dim@x+\dim@sep-\dim@overlineII}{\dim@y}}
		\fi
		% Draw an extra line if node text lies outside of the
		% dimension.
		\pgfmathparse{abs(.5\dim@y)<abs(\dim@text@translation)}
		\if\pgfmathresult1
			% Dimension label outside.
			\ifdim0pt<\dim@text@translation
				% Text lies on right.
				\pgfpathmoveto{\pgfpointanchor{\dim@text@nodename}{south east}}
				\pgfpathlineto{\pgfpoint{\dim@x+\dim@sep}{\dim@y}}
			\else
				% Text lies on left.
				\pgfpathmoveto{\pgfpointanchor{\dim@text@nodename}{south west}}
				\pgfpathlineto{\pgfpoint{\dim@x+\dim@sep}{0pt}}
			\fi
		\fi
		% Draw the dimension segments with arrows.
		\pgfsetarrowsstart{\dim@arrow@type}
		\pgfsetarrowsend{\dim@arrow@type}
		\pgfpathmoveto{\pgfpoint{\dim@x+\dim@sep}{0pt}}
		\pgfpathlineto{\pgfpoint{\dim@x+\dim@sep}{\dim@y}}
		% Finally put all the stuff on paper.
		\pgfusepath{stroke}
	\fi
\else
	%
	% FIRST POINT LOWER
	%
	\ifnum90>\pgfdecoratedangle
		% First quadrant
		% Node drawing. It can be translate along the path with
		% text translation option.
		{
 		\pgftransformshift{\pgfpoint{\dim@x+\dim@sep}{\dim@y/2-\dim@text@translation}}
		% In this case, the node must e rotated.
		\pgftransformrotate{-90}
		% Capturing the TikZ picture font size
		\dim@text@font
		\pgfnode{rectangle}{south}{\dim@text}{\dim@text@nodename}{\pgfusepath{discard}}
		}
		% Drawing the connection segments.
		\pgfpathmoveto{\pgfpoint{\dim@x+\dim@sep+\dim@overlineI}{0pt}}
		% Check if the first dimension connection must be drawn.
		\if@dim@connection
			\pgfpathlineto{\pgfpoint{0pt}{0pt}}
		\else
			\pgfpathlineto{\pgfpoint{\dim@x+\dim@sep-\dim@overlineI}{0pt}}
		\fi
		\pgfpathmoveto{\pgfpoint{\dim@x+\dim@sep+\dim@overlineII}{\dim@y}}
		% Check if the second dimension connection must be drawn.
		\if@dim@connection
			\pgfpathlineto{\pgfpoint{\dim@x}{\dim@y}}
		\else
			\pgfpathlineto{\pgfpoint{\dim@x+\dim@sep-\dim@overlineII}{\dim@y}}
		\fi
		% Draw an extra line if node text lies outside of the
		% dimension.
		\pgfmathparse{abs(.5\dim@y)<abs(\dim@text@translation)}
		\if\pgfmathresult1
			% Dimension label outside.
			\ifdim0pt<\dim@text@translation
				% Text lies on right.
				\pgfpathmoveto{\pgfpointanchor{\dim@text@nodename}{south east}}
				\pgfpathlineto{\pgfpoint{\dim@x+\dim@sep}{\dim@y}}
			\else
				% Text lies on left.
				\pgfpathmoveto{\pgfpointanchor{\dim@text@nodename}{south west}}
				\pgfpathlineto{\pgfpoint{\dim@x+\dim@sep}{0pt}}
			\fi
		\fi
		% Draw the dimension segments with arrows.
		\pgfsetarrowsstart{\dim@arrow@type}
		\pgfsetarrowsend{\dim@arrow@type}
		\pgfpathmoveto{\pgfpoint{\dim@x+\dim@sep}{0pt}}
		\pgfpathlineto{\pgfpoint{\dim@x+\dim@sep}{\dim@y}}
		% Finally put all the stuff on paper. 
		\pgfusepath{stroke}
	\else 
		% Second quadrant
		% Node drawing. It can be translate along the path with
		% text translation option.
		{
 		\pgftransformshift{\pgfpoint{\dim@sep}{\dim@y/2-\dim@text@translation}}
		% In this case, the node must e rotated.
		\pgftransformrotate{-90}
		% Capturing the TikZ picture font size
		\dim@text@font
		\pgfnode{rectangle}{south}{\dim@text}{\dim@text@nodename}{\pgfusepath{discard}}
		}
		% Drawing the connection segments.
		\pgfpathmoveto{\pgfpoint{\dim@sep+\dim@overlineI}{0pt}}
		% Check if the first dimension connection must be drawn.
		\if@dim@connection
			\pgfpathlineto{\pgfpoint{0pt}{0pt}}
		\else
			\pgfpathlineto{\pgfpoint{\dim@sep-\dim@overlineI}{0pt}}
		\fi
		\pgfpathmoveto{\pgfpoint{\dim@sep+\dim@overlineII}{\dim@y}}
		% Check if the second dimension connection must be drawn.
		\if@dim@connection
			\pgfpathlineto{\pgfpoint{\dim@x}{\dim@y}}
		\else
			\pgfpathlineto{\pgfpoint{\dim@sep-\dim@overlineII}{\dim@y}}
		\fi
		% Draw an extra line if node text lies outside of the
		% dimension.
		\pgfmathparse{abs(.5\dim@y)<abs(\dim@text@translation)}
		\if\pgfmathresult1
			% Dimension label outside.
			\ifdim0pt<\dim@text@translation
				% Text lies on right.
				\pgfpathmoveto{\pgfpointanchor{\dim@text@nodename}{south east}}
				\pgfpathlineto{\pgfpoint{\dim@sep}{0pt}}
			\else
				% Text lies on left.
				\pgfpathmoveto{\pgfpointanchor{\dim@text@nodename}{south west}}
				\pgfpathlineto{\pgfpoint{\dim@sep}{\dim@y}}
			\fi
		\fi
		% Draw the dimension segments with arrows.
		\pgfsetarrowsstart{\dim@arrow@type}
		\pgfsetarrowsend{\dim@arrow@type}
		\pgfpathmoveto{\pgfpoint{\dim@sep}{0pt}}
		\pgfpathlineto{\pgfpoint{\dim@sep}{\dim@y}}
		% Finally put all the stuff on paper. 
		\pgfusepath{stroke}
	\fi
\fi
}}

%-> Along the dimension
\pgfdeclaredecoration{dim}{final}{
\state{final}{%
% Check if the dimension inserted has unit
\tikz@checkunit{\pgfkeysvalueof{/pgf/decoration/distance}}
\iftikz@isdimension
	\pgfmathsetlength{\dim@sep}{\pgfkeysvalueof{/pgf/decoration/distance}}
\else
	\pgfmathsetlength{\dim@sep}{\pgfkeysvalueof{/pgf/decoration/distance}*\pgf@yy}
\fi
% Overline value if function of the sep sign.
\ifdim0pt>\dim@sep
	\pgfmathsetlength{\dim@overline}{-\pgfkeysvalueof{/pgf/decoration/overline}}
\else
	\pgfmathsetlength{\dim@overline}{\pgfkeysvalueof{/pgf/decoration/overline}}
\fi
\pgfmathsetlength{\dim@text@translation}{\pgfkeysvalueof{/pgf/decoration/text translation}}
% Setting the text to be inserted into the dimension node.
\def\dim@text{\pgfkeysvalueof{/pgf/decoration/text}}
{
\pgftransformshift{\pgfpoint{\pgfdecoratedpathlength/2+\dim@text@translation}{\dim@sep}}
% Capturing the TikZ picture font size
\dim@text@font
\pgfnode{rectangle}{south}{\dim@text}{\dim@text@nodename}{\pgfusepath{discard}}
}
\pgfpathmoveto{\pgfpoint{0pt}{\dim@sep+\dim@overline}}
% Check if the first connection line mus b drawn
\if@dim@connection
	\pgfpathlineto{\pgfpoint{0pt}{0pt}}
\else
	\pgfpathlineto{\pgfpoint{0pt}{\dim@sep-\dim@overline}}
\fi
\pgfpathmoveto{\pgfpoint{(\pgfdecoratedpathlength}{\dim@sep+\dim@overline}}
% Check if the first connection line mus b drawn
\if@dim@connection
	\pgfpathlineto{\pgfpoint{\pgfdecoratedpathlength}{0pt}}
\else
	\pgfpathlineto{\pgfpoint{(\pgfdecoratedpathlength}{\dim@sep-\dim@overline}}
\fi
% Draw an extra line if node text lies outside of the
% dimension.
\pgfmathparse{abs(.5*\pgfdecoratedpathlength)<abs(\dim@text@translation)}
\if\pgfmathresult1
	% Dimension label outside.
	\ifdim0pt<\dim@text@translation
		% Text lies on right.
		\pgfpathmoveto{\pgfpointanchor{\dim@text@nodename}{south east}}
		\pgfpathlineto{\pgfpoint{\pgfdecoratedpathlength}{\dim@sep}}
	\else
		% Text lies on left.
		\pgfpathmoveto{\pgfpointanchor{\dim@text@nodename}{south west}}
		\pgfpathlineto{\pgfpoint{0pt}{\dim@sep}}
	\fi
\fi
\pgfsetarrowsstart{\dim@arrow@type}
\pgfsetarrowsend{\dim@arrow@type}
\pgfpathmoveto{\pgfpoint{0pt}{\dim@sep}}
\pgfpathlineto{\pgfpoint{\pgfdecoratedpathlength}{\dim@sep}} 
\pgfusepath{stroke}
\pgfpathmoveto{\pgfpoint{0pt}{0pt}}
}}

%-> Initial values
\pgfkeys{/pgf/decoration/.cd,
         distance/.initial         = 10pt,
         overline/.initial         = 1mm,
         text/.initial             = {},
         text translation/.initial = 0pt,
         text node name/.store in  = \dim@text@nodename,
         text node name            = dim_text,
         >/.store in               = \dim@arrow@type,
         >                         = latex,
         connection/.is if         = @dim@connection,
         connection                = false,
         font/.store in            = \dim@text@font,
         font                      = \tikz@textfont,
}
\makeatother

%%%%%%%%%%%%%%%%%%%%%%%%%%%%%%%%%%%%%%%%%%
% COMMANDS TO CALL DIMENSION DECORATIONS %
%%%%%%%%%%%%%%%%%%%%%%%%%%%%%%%%%%%%%%%%%%
\newcommand\Hdimension[1][]{\path[decorate,decoration={Hdim,#1}]}
\newcommand\Vdimension[1][]{\path[decorate,decoration={Vdim,#1}]}
\newcommand\dimension[1][]{\path[decorate,decoration={dim,#1}]}

\begin{document}
\begin{tikzpicture}[x         = 1mm,
                    y         = 1mm,
                    >         = latex,
                    line join = round,
                    font      = \small]

%%%%%%%%%%%%%%%%%%
% NOZZLE DRAWING %
%%%%%%%%%%%%%%%%%%
%-> Origin definition
\coordinate (o) at (0,0);

% Name: Supersonic Nozzle
% Description: Used to speed up liquids by reducing the volume
% Functions: 
%	F1. Mounting : The Nozzle can be mounted
%	F2. Presure increase : The presure of the liquid gets increased
%	F3. Velocity increase:  The velocity of the liquid gets increased
% Principals:
%	W1. conservation of energy : the energy stays always the same
%	W2. Conservation of mass : the mass stays always the same
%	W3. Use bernoulli equation to calculate the presure levels at different points
%	W4. Continuity equation to calculate the velocities
% Parts:
%	P1. Nozzle Inlet : The inlet of the nozzle
%	P2. Nozzle Outlet : The outlet of the nozzle
%	P2.1 Throat : The slope of the diverging side of the de laval nozzle
%	P3. Flange : A flange to mount the nozzle
%	P4. Retaining groove : Enables connection to the inlet
% Visual:
%	A side view of a supersonic nozzle, cut in the middle
%	V1. Hatched nozzle body
%	V2. Dimensional drawings
%	V3. A symetrie line in the center of the nozzle
%	V4. Lines for mounting screws
%	V5. Symetric axis of the screw holes


%Definition of the nozzle
% Parameter:
\newcommand{\lengthBody}{100} % Length of the nozzle from flange to outlet 
\newcommand{\bodyTipChamfer}{10} % Horizontal length of the nozzle tip

% Nozzle inlet
\newcommand{\inletDiameter}{23} % Defines the inlet diameter of the supersonic nozzle


% Nozzle outlet
\newcommand{\outletDiameter}{10} % Defines the outlet diameter of the nozzle.
\newcommand{\outletLength}{43.8} % Defines the outlet diameter of the nozzle.

% Flange variables
\newcommand{\flangeDiameter}{80} % Defines the diameter of the flange
\newcommand{\flangeScrewHoleDiameter}{8} % Diameter of the screw holes
\newcommand{\flangeScrewHoleCount}{4} % Count of holes inside the flange
\newcommand{\flangeHoleDiameter}{60} % Defines the diameter of screw hole positions
\newcommand{\flangeRimSize}{6} % The size of the flange rim
\newcommand{\flangeInnerDiameter}{11} % Size of the inner ring of the flange
\newcommand{\flangeWallthickness}{10} % Wallthickness of the flange

% Retaining groove variables
\newcommand{\rGrooveThickness}{5} % Wallthicknes of the Retaining groove
\newcommand{\rGrooveLength}{20} % The length of the retaining groove
\newcommand{\rGrooveNotchDepth}{1.575} % The depth of the nodge between retain groove and flange
\newcommand{\rGrooveNotchLength}{10.0} % The length of the nodge
\newcommand{\rGrooveOutDiameter}{\fpeval{\rGrooveNotchDepth +\inletDiameter}} % The calculated out diameter of the retaining groove

%Throat variables
\newcommand{\throatRadius}{2.5} % The radius of the throat
\newcommand{\throatDiameter}{\fpeval{\throatRadius*2}} % Diameter by radius
\newcommand{\throatOffset}{\fpeval{\throatRadius}} % Helper,Offset for drawing purpose



%Throat taper
\newcommand{\throatTaperRadius}{7} % Radius of the taper arcs between inlet and outlet

% Constants
\newcommand{\outletRadius}{\fpeval{\outletDiameter/2}} % Calculate the radius of the outlet
\newcommand{\inletRadius}{\fpeval{\inletDiameter/2}} % Calculate the radius of the inlet
\newcommand{\inletLength}{\fpeval{\rGrooveLength+\rGrooveNotchLength+\flangeWallthickness+\lengthBody+\bodyTipChamfer -\outletLength-\throatTaperRadius-\throatTaperRadius}} % Calculate the length of inlet


%-> Nozzle
% Symmetric characteristic is used in a
% foreach command where the cycle is made
% for upper part (up) with positive sign (+)
% and lower part (down) with negative sign (-)
\foreach \pos/\sign in {up/+,down/-}{

%%-> Points definitions
\coordinate (A\pos)        at ($(o)+(0,\sign\inletRadius)$); % Start inlet chamber left corner, Retaining groove left bottom corner
\coordinate (B\pos)        at ($(A\pos)+(0,\sign\rGrooveThickness)$); % Top left corner Retaining groove
\coordinate (C\pos)        at ($(B\pos)+(\rGrooveLength,0)$);  %Top right corner retaining groove
\coordinate (D\pos)        at ($(C\pos)+(0,-\sign\rGrooveNotchDepth)$); % Right bottom corner retaining groove
\coordinate (E\pos)        at ($(D\pos)+(\rGrooveNotchLength,0)$); %Bottom right corner retaining groove, bottom left corner inner flange
\coordinate (F\pos)        at ($(E\pos)+(0,\sign\flangeInnerDiameter)$); %Top left corner inner flange, bottom left corner screw hole
\coordinate (G\pos)        at ($(F\pos)+(0,\sign\flangeScrewHoleDiameter)$); %Screw hole left top corner, Flange rim bottom left corner
\coordinate (H\pos)        at ($(G\pos)+(0,\sign\flangeRimSize)$); %Flange rim left top corner
\coordinate (I\pos)        at ($(H\pos)+(\flangeWallthickness,0)$); % Right top corner of the flange rim
\coordinate (L\pos)        at (G\pos-|I\pos); % Flange rim bottom left corner, Screw hole top right corner
\coordinate (M\pos)        at (F\pos-|I\pos); % Screw hole right bottom corner, Flange inner top right corner
\coordinate (N\pos)        at (E\pos-|I\pos); % Flange inner bottom right corner, Nozzle body top left corner 
\coordinate (O\pos)        at ($(N\pos)+(\lengthBody,0)$); % Nozzle body right corner, Upper tip corner
\coordinate (P\pos)        at ($(o)+(\rGrooveLength+\rGrooveNotchLength+\flangeWallthickness+\lengthBody+\bodyTipChamfer,\sign\outletRadius)$);  % End of outlet, tip corner
\coordinate (throat_\pos)  at ($(o)+(\rGrooveLength+\rGrooveNotchLength+\flangeWallthickness+\lengthBody+\bodyTipChamfer -\outletLength,\sign\throatOffset)$); % Throat position
\coordinate (IN_\pos)      at ($(o)+((\inletLength,\sign\inletRadius)$); % End inlet chamber, point next to taper
\coordinate (center1_\pos) at ($(o)+(\inletLength,\sign\inletRadius-\sign\throatTaperRadius)$); % Taper top Arc radius position
\coordinate (center2_\pos) at ($(throat_\pos)+(0,\sign\throatTaperRadius)$);  % Taper bottom Arc radius position 


%%-> Draw nozzle body
%% Perspektive: Side-view
\draw[fill,
      pattern    = section, % Use hatching
      line width = 1.1pt] % Outer line width of the nozzle
\ifnum\sign1>0 % Make the drawing symetrical along the Y axis
        (center1_\pos)++(45:\throatTaperRadius)arc(45:90:\throatTaperRadius)-- % Draw the top outer arc of the curved diameter reduction
\else
        (center1_\pos)++(-45:\throatTaperRadius)arc(315:270:\throatTaperRadius)--% Draw the bottom outer arc of the curved diameter reduction
\fi
(A\pos)-- % Inner wall inlet chamber, Line from end of taper arc to the inlet side 
(B\pos)-- % Outer left wall Inlet and the Retaining groove
(C\pos)-- % Retaining groove upper wall
(D\pos)-- % Retaining groove right wall, left wall of the nodge
(E\pos)-- % Nodge bottom
(F\pos)-- % Left wall of the Flange inner
(G\pos)-- % Left line of the screw hole
(H\pos)-- % Left wall of the Flange rim
(I\pos)-- % Upper wall of the Flange rim
(L\pos)-- % Right wall of the Flange rim
(M\pos)-- % Right line of the screw hole
(N\pos)-- % Right wall of the Flange inner
(O\pos)-- % Upper wall of the Nozzle body
(P\pos)-- % Nozzle tip diagonal line
(throat_\pos) % Inner wall of the outler
\ifnum\sign1>0 
      arc(270:225:\throatTaperRadius)  % Draw the top inner arc of the throat taper
\else
     arc(90:135:\throatTaperRadius) % Draw the bottom inner arc of the throat taper
\fi
--cycle;

%%-> Draw screw holes on the flange as overlay
\draw[fill       = white,
      line width = 1.1pt] (G\pos)rectangle(M\pos); % Draw the screwholes as white rectangle from the bottom left top the right top corner of the screw hole

%%-> Draw symmetric axis on flange screw holes
\draw[dash pattern = on 3pt off 5pt on 6pt off 5pt,
      line width   = 1pt] ($(F\pos)!.5!(G\pos)-(5,0)$)--
                          ($(L\pos)!.5!(M\pos)+(5,0)$);

%%-> screw drawing
\draw[dashed] (D\pos)--(D\pos-|A\pos); % Screw inside Retaining groove

}

%%-> Nozzle input and output closure
\draw[line width = 1.1pt](Adown)--(Aup)
                         (Pdown)--(Pup); % Closure lines on the left on right side of the groove

%%-> Nozzle symmetry line
\draw[dash pattern = on 3pt off 5pt on 6pt off 5pt,
      line width   = 1pt] ($(Adown)!.5!(Aup)-(5,0)$)--($(Pdown)!.5!(Pup)+(5,0)$); % Symetrie line in the center of the nozzle.

%%%%%%%%%%%%%%
% DIMENSIONS %
%%%%%%%%%%%%%%
% Macro to see the dimension
% inserted. For debug.
\newif\ifdimension
\dimensionfalse
%\dimensiontrue % if true, you will see the dimension number (%x) on the draw
\newcount\Ndim=0
\def\SeeDim#1{\ifdimension\global\advance\Ndim by 1 \the\Ndim\else#1\fi}
%-> 1
\Hdimension[text      = \SeeDim{\flangeWallthickness},
            distance  = 3]  (Hup) --  (Iup);

%-> 2
\Hdimension[text     = \SeeDim{\rGrooveNotchLength},
            distance = 3] (Cup)--(Hup);

%-> 3
\Hdimension[text     = \SeeDim{\rGrooveLength},
            distance = 26.425] (Bup)--(Cup);

%-> 4
\Hdimension[text     = \SeeDim{\lengthBody},
            distance = 3] (Iup)--(Oup);

%-> 5
\Hdimension[text     = \SeeDim{\bodyTipChamfer},
            distance = 28] (Oup)--(Pup);

%-> 6
\dimension[text     = \SeeDim{\diameter30},
           distance = -25] (Odown)--(Oup);

%-> 7
\Hdimension[text     = \SeeDim{\outletLength},
            distance = -49.5] (throat_down)--(Odown);

%-> 8
\Hdimension[text     = \SeeDim{\inletLength},
            distance = -40.5] (IN_down)--(Bdown);

%-> 9
\dimension[text             = \SeeDim{\diameter\throatDiameter},
           text translation = -7mm,
           distance         = 20] (throat_down)--(throat_up);

%-> 10


\dimension[text     = \SeeDim{\diameter \outletDiameter}, % Outlet Diameter Text
           distance = -8] (Pdown)--(Pup);

%-> 11 (by hand)
\draw[->] (center1_up)--++(45:7)node[sloped,
                                     above     = .4mm,
                                     pos       = .3,
                                     inner sep = 0.5pt]{\SeeDim{R7}};

%-> 12 (by hand)
\draw[->] (center2_up)--++(225:7)node[above     = .4mm,
                                      pos       = .4,
                                      inner sep = 0.5pt,
                                      rotate    = 45,
                                      fill      = white]{\SeeDim{R7}};

%-> 13 (by hand)
\coordinate (raccordo_up)   at ($(center1_up)+(45:7)$);
\coordinate (raccordo_down) at ($(center1_down)+(-45:7)$);
\draw (raccordo_up)--++(135:20);
\draw (raccordo_down)--++(225:20);
\path[name path=C1](raccordo_up)--++(-45:20);
\path[name path=C2](raccordo_down)--++(45:20);
\path[name intersections={of=C1 and C2}];
\coordinate (C90) at (intersection-1);
\draw[<->] ($(C90)+(135:30)$)arc[start angle = 135,
                                 delta angle = 90,
                                 radius      = 30];
\def\angle{35}
\node[rotate = \angle+45,
      anchor = south] at ($(C90)+(135+\angle:30)$){\SeeDim{\ang{90}}};

%-> 14
\dimension[text             = \SeeDim{\diameter\flangeHoleDiameter},
           distance         = -5,
           >                = angle 45,
           text translation = .5cm] ($(Ldown)!.5!(Mdown)$)--($(Lup)!.5!(Mup)$);

%-> 15
\dimension[text             = \SeeDim{\diameter\flangeDiameter},
           distance         = -10,
           text translation = .5cm] (Idown)--(Iup);

%-> 16
\dimension[text             = \SeeDim{\diameter\flangeScrewHoleDiameter ($\times$\flangeScrewHoleCount)},
           distance         = 10,
           text translation = -12mm] (Gdown)--(Fdown);

%-> 17
\dimension[text     = \SeeDim{\diameter\rGrooveOutDiameter},
           distance = 10] (Bdown)--(Bup);

%-> 18
\dimension[text     = \SeeDim{\diameter\inletDiameter},
           distance = 5] (Adown)--(Aup);

% Dimensions scale
\node[anchor = north west] at (current bounding box.south west)
{All dimensions are in millimeters};
\end{tikzpicture}
\end{document}
